\begin{proposition}[Closedness Is Necessary in the Nested Interval Theorem]
\label{prop:closedness-is-necessary-in-nested-interval-theorem}
There exists a sequence $(I_n)_{n\in\mathbb{N}}$ of nonempty bounded intervals
in $\mathbb{R}$, not all of which are closed, such that
\[
  (\forall n\in\mathbb{N})(I_{n+1}\subseteq I_n),
\]
but
\[
  \bigcap_{n=1}^{\infty}I_n=\varnothing.
\]
For example,
\[
  I_n=\left(0,\frac{1}{n}\right].
\]
\hyperref[prf:closedness-is-necessary-in-nested-interval-theorem]{Go to proof.}
\end{proposition}
\end{propositionbox}

\begin{remark*}[Standard quantified statement]
\[
\exists I_n\;\exists\mathcal{I}\;
\bigl(
\mathcal{I}=\mathsf{IndexedFamily}(I_n,\mathbb{N},\mathbb{R})
\land
\operatorname{MembersSatisfy}(\mathcal{I},\mathsf{NonemptyBoundedInterval},\mathbb{R})
\land
\operatorname{Nested}(\mathcal{I},\subseteq)
\land
\neg\operatorname{IsClosedInterval}(I_n)
\land
\bigcap_{n=1}^{\infty}I_n=\varnothing
\bigr).
\]
\end{remark*}

\begin{remark*}[Predicate reading]
\[
\mathcal{I}=\mathsf{IndexedFamily}(I_n,\mathbb{N},\mathbb{R}),\qquad C=\bigcap_{n=1}^{\infty}I_n,\qquad
\neg\operatorname{IsClosedInterval}(I_n,\mathbb{R})
\quad\text{can make}\quad
\operatorname{IsNonempty}(C)
\text{ fail.}
\]
\end{remark*}

\begin{remark*}[Negated quantified statement]
\[
\forall I_n\;\forall\mathcal{I}\;
\left[
\mathcal{I}=\mathsf{IndexedFamily}(I_n,\mathbb{N},\mathbb{R})
\land
\operatorname{MembersSatisfy}(\mathcal{I},\mathsf{NonemptyBoundedInterval},\mathbb{R})
\land
\operatorname{Nested}(\mathcal{I},\subseteq)
\Longrightarrow
\bigcap_{n=1}^{\infty}I_n\ne\varnothing
\right].
\]
\end{remark*}

\begin{remark*}[Negation predicate reading]
\[
C=\bigcap_{n=1}^{\infty}I_n,\qquad F=(I_n)_{n\in\mathbb{N}},\qquad
P_{\mathbb{R}}=\mathsf{OrderedSet}(\mathbb{R},\leq),
\]
\[
C=\varnothing
\Longrightarrow
\neg\operatorname{NestedClosedBoundedIntervals}(F,P_{\mathbb{R}}).
\]
\end{remark*}

\begin{remark*}[Failure modes]
\begin{description}
  \item[Exposition.]
  Failure is witnessed by the displayed negated or contrapositive condition.

  \item[\operatorname{MembersSatisfy} / \operatorname{Nested}.]
  \[
\forall I_n\;\forall\mathcal{I}\;
\left[
\mathcal{I}=\mathsf{IndexedFamily}(I_n,\mathbb{N},\mathbb{R})
\land
\operatorname{MembersSatisfy}(\mathcal{I},\mathsf{NonemptyBoundedInterval},\mathbb{R})
\land
\operatorname{Nested}(\mathcal{I},\subseteq)
\Longrightarrow
\bigcap_{n=1}^{\infty}I_n\ne\varnothing
\right].
\]
\[
C=\bigcap_{n=1}^{\infty}I_n,\qquad F=(I_n)_{n\in\mathbb{N}},\qquad
P_{\mathbb{R}}=\mathsf{OrderedSet}(\mathbb{R},\leq),
\]
\end{description}
\end{remark*}

\begin{remark*}[Interpretation]
Dropping closedness can remove the endpoint that would otherwise remain in
the limiting intersection.
\end{remark*}

\begin{dependencies}
\hyperref[def:open-interval]{Open Interval};
\hyperref[def:bounded-subset-of-r]{Bounded Subset of the Real Line}.
\end{dependencies}

